\section{Špecifikácia požiadaviek a~výber hardvéru}

Pred samotným návrhom a~implementáciou riešenia bolo
potrebné špecifikovať, aké funkcie má aplikácia plniť
a~aké požiadavky musí spĺňať hardvérové zariadenie,
na ktorom bude bežať časť aplikácie. V~tejto kapitole
postupne definujeme funkcionálne a~nefunkcionálne
požiadavky na aplikáciu, z~nich odvodené hardvérové
požiadavky a~následne predstavíme zariadenie zvolené
pre praktickú implementáciu.

\subsection{Funkcionálne požiadavky}
\label{subsec:funkcionalne-poziadavky}

Funkcionálne požiadavky definujú, čo má aplikácia robiť
z~pohľadu používateľa. Na základe analýzy pravidiel
zápasenia (kapitola~\ref{chapter:zapasenie}) a~potrieb
rozhodcu sme stanovili nasledujúce funkcionálne
požiadavky:

\begin{itemize}
  \item \textbf{Rozpoznávanie gest rozhodcu} -- aplikácia
        musí rozpoznať osem gest zodpovedajúcich akciám
        rozhodcu počas zápasu: bodovanie (1, 2 a~4~body
        pre červeného aj modrého zápasníka), napomenutie
        červeného a~modrého zápasníka, pasivitu červeného
        a~modrého zápasníka a~lopatkové víťazstvo TOUCHE.

  \item \textbf{Aktualizácia stavu zápasu} -- po detekcii
        gesta musí aplikácia automaticky aktualizovať
        aktuálne skóre a~záznamy o~napomenutiach
        zápasníkov.

  \item \textbf{Zobrazovanie stavu zápasu} -- aplikácia
        musí na mobilnom telefóne priebežne zobrazovať
        aktuálne skóre oboch zápasníkov, počet
        napomenutí a~prípadné lopatkové víťazstvo.

  \item \textbf{Spätná väzba pre rozhodcu} -- aplikácia
        musí poskytovať rozhodcovi okamžité potvrdenie
        o~detekcii gesta priamo na hodinkách
        (napríklad vibračnou odozvou) bez nutnosti
        pozerania na obrazovku.

  \item \textbf{Záznam udalostí} -- aplikácia musí
        zaznamenať každú detegovanú udalosť s~časovou
        pečiatkou pre potreby spätnej kontroly priebehu
        zápasu.
\end{itemize}

\subsection{Nefunkcionálne požiadavky}
\label{subsec:nefunkcionalne-poziadavky}

Nefunkcionálne požiadavky určujú, ako má aplikácia
plniť svoje funkcionálne úlohy -- definujú kvalitatívne
vlastnosti ako rýchlosť, spoľahlivosť či použiteľnosť:

\begin{itemize}
  \item \textbf{Rýchlosť odozvy} -- detekcia gesta
        a~následná aktualizácia stavu zápasu musia
        prebehnúť v~rádoch sekúnd od dokončenia
        pohybu rozhodcu, aby nedošlo k~oneskoreniu
        narúšajúcemu priebeh zápasu.

  \item \textbf{Spoľahlivosť detekcie} -- aplikácia musí
        dosahovať dostatočne vysokú úspešnosť rozpoznávania
        gest, aby bola prakticky použiteľná. Konkrétne
        výsledky testovania uvádzame v~kapitole
        \ref{chapter:testovanie}.

  \item \textbf{Odolnosť voči falošným detekciám} --
        aplikácia nesmie nesprávne detegovať gestá pri
        bežných pohyboch ruky, ktoré nesúvisia
        so~signalizáciou rozhodcu.

  \item \textbf{Použiteľnosť} -- aplikácia musí byť
        použiteľná bez zložitej konfigurácie pred zápasom
        a~bez nutnosti programátorských znalostí zo strany
        rozhodcu.

  \item \textbf{Energetická efektívnosť} -- hodinky aj
        mobilný telefón musia byť schopné zvládnuť aspoň
        jeden zápas (cca~7~minút) bez nutnosti nabíjania.
\end{itemize}

\subsection{Hardvérové a~platformové požiadavky}
\label{subsec:hw-poziadavky}

Keďže základom aplikácie je vyhodnocovanie pohybových
gest rozhodcu priamo zo zápästia, nie každé smart
nositeľné zariadenie je pre tento účel vhodné.
Z~funkcionálnych a~nefunkcionálnych požiadaviek
vyplynuli nasledujúce požiadavky na inteligentné
hodinky:

\begin{itemize}
  \item \textbf{Prístup k~akcelerometru a~gyroskopu
        cez API} -- zariadenie musí umožňovať
        programátorský prístup k~surovým dátam
        zo senzorov pohybu, nie len ich spracované
        výstupy (napr. počet krokov). Bez tejto možnosti
        nie je možné vyhodnocovať vlastné gestá.

  \item \textbf{Podpora vlastných aplikácií} -- hodinky
        musia disponovať otvoreným vývojovým prostredím
        umožňujúcim inštaláciu aplikácií tretích strán.
        Uzavreté platformy bez \Gls{sdk} či developer
        \Gls{api} vývojárom neumožňujú nasadiť vlastný
        kód.

  \item \textbf{Komunikácia s~mobilným zariadením} --
        aplikácia vyžaduje prenos dát z~hodiniek
        do~Androidu v~reálnom čase. Hodinky preto musia
        podporovať obojsmernú Bluetooth komunikáciu
        prístupnú z~kódu aplikácie.

  \item \textbf{Dostatočná frekvencia snímania senzorov} --
        na zachytenie krátkych a~rýchlych pohybov
        charakteristických pre rozhodcovské gestá je
        potrebná dostatočne vysoká vzorkovacia frekvencia.
        Príliš nízka frekvencia by viedla k~strate
        dôležitých pohybových informácií.

  \item \textbf{Výdrž batérie} -- zariadenie musí byť
        schopné pracovať počas celého trvania zápasu
        bez nutnosti nabíjania, čo kladie požiadavku
        na energetickú efektívnosť operačného systému
        aj aplikácie samotnej.
\end{itemize}

\subsection{Amazfit T-Rex 3}

Po splnení všetkých hardvérových a~platformových
požiadaviek sme pre praktickú implementáciu zvolili
hodinky Amazfit T-Rex 3. Z~pohľadu hardvéru a~softvéru
ponúkajú parametre, ktoré sú pre našu aplikáciu
dostatočné a~zároveň umožňujú stabilný vývoj
v~prostredí Zepp OS 4.

\begin{itemize}
  \item \textbf{Procesor:} ZPS3044
  \item \textbf{Operačná pamäť:} 64 MB \Gls{ram}
  \item \textbf{Úložisko:} 32 GB
  \item \textbf{Displej:} 1,5" \Gls{amoled}, 480x480 px, 2000 nits
  \item \textbf{Senzory:} 6-osový \Gls{imu} (3-osový
        akcelerometer + 3-osový gyroskop), barometer,
        kompas, optický senzor srdcovej frekvencie
  \item \textbf{Konektivita:} Bluetooth 5.3 \Gls{ble}
  \item \textbf{Operačný systém:} Zepp OS 4
  \item \textbf{Výdrž batérie:} až 27 dní pri bežnom
        používaní
\end{itemize} \cite{trex3_specs}
